\documentclass[11pt]{article}

% Pacotes
\RequirePackage{luatex85}
\usepackage{amsmath}
\usepackage{amssymb}
\usepackage{graphicx}
\usepackage[colorlinks=true, allcolors=blue]{hyperref}
\usepackage{amsthm}
\usepackage{framed, color,xcolor}
\usepackage{stmaryrd}
\usepackage{rotating}           % Usado para rotacionar o texto
\usepackage[all,knot,arc,import,poly]{xy}   % Pacote para desenhos gráficos
% Este pacote pode conflitar com outros pacotes gráficos como o ``pictex''
% Então é necessário usar apenas um dos pacotes conflitantes
\usepackage{tikz, tikz-cd}
\usepackage{epigraph}
\usetikzlibrary{babel}
\usetikzlibrary{positioning}
\usetikzlibrary{calc}
\usepackage[stix2]{newtxmath}
\usepackage[usenames,dvipsnames]{xcolor}
\usepackage{mathtools, nccmath}
\usepackage{amssymb, mathrsfs}
\usepackage[numbers]{natbib}
\usepackage{fontspec}
\usepackage{enumitem} % pra substituir o símbolo do itemize por alguma coisa legal, tipo setinhas:  \begin{itemize}[label=\ding{212}]
\usepackage{pifont}
\usepackage{stmaryrd} % for \arrownot (notação de adjacência em grafo)
\usepackage{quiver} %pacote do quiver pra desenhar diagramas mais legais

%Fontes e Línguas
\usepackage[portuguese]{babel}
\usepackage{eufrak}
\setmainfont{DarwinSerif-Regular.otf}
\newfontfamily\fakebold[FakeBold=1.75]{DarwinSerif-Regular.otf}
\newfontfamily\fakeitalic[FakeSlant=0.2]{DarwinSerif-Regular.otf}
\renewcommand{\textbf}[1]{{\fakebold#1}}
\renewcommand{\textit}[1]{{\fakeitalic#1}}
\renewcommand{\bfseries}{\fakebold}
\renewcommand{\itshape}{\fakeitalic}



%Cores


%Comandinhos que gosto
\usepackage{mathtools}
\newcommand{\defeq}{\vcentcolon=}
\newcommand{\eqdef}{=\vcentcolon}
\newcommand{\bigslant}[2]{{\raisebox{.2em}{$#1$}\left/\raisebox{-.2em}{$#2$}\right.}}
\newcommand{\logo}{\Rightarrow}
\newcommand{\pulinho}{\vspace{.5cm}}
\newcommand{\edge}{%
  \mathrel{-}% minus as a relation
  \joinrel\joinrel % some backup
  \mathrel{-}% minus as a relation
}
\newcommand{\notedge}{%
  \mathrel{\mkern2mu}% a small advancement
  \arrownot % a short slash
  \mathrel{\mkern-2mu}% compensate
  \edge
}

\newcommand{\Aut}{\text{Aut}}
\newcommand{\Cay}{\text{Cay}}

\newcommand{\luafacil}{
oi
}

\newcommand{\luamedio}{
oi
}

\let\varphi\phi

%Caixinhas coloridas

\theoremstyle{definition}
\newtheorem{theorem}{Teorema}
\newenvironment{teorema}{
\definecolor{shadecolor}{HTML}{f59aa3}
\begin{shaded}
\begin{theorem}
}
{
\end{theorem}
\end{shaded}
}

\theoremstyle{definition}
\newtheorem*{theorem*}{Teorema}
\newenvironment{teorema*}{
\definecolor{shadecolor}{HTML}{f59aa3}
\begin{shaded}
\begin{theorem*}
}
{
\end{theorem*}
\end{shaded}
}

\theoremstyle{definition}
\newtheorem{prop}{Proposição}
\newenvironment{proposicao}{
\definecolor{shadecolor}{HTML}{FFDBDF}
\begin{shaded}
\begin{prop}
}
{
\end{prop}
\end{shaded}
}

\theoremstyle{definition}
\newtheorem*{prop*}{Proposição}
\newenvironment{proposicao*}{
\definecolor{shadecolor}{HTML}{FFDBDF}
\begin{shaded}
\begin{prop*}
}
{
\end{prop*}
\end{shaded}
}

\theoremstyle{definition}
\newtheorem{cor}{Corolário}
\newenvironment{corolario}{
\definecolor{shadecolor}{HTML}{f59aa3}
\begin{shaded}
\begin{cor}
}
{
\end{cor}
\end{shaded}
}

\theoremstyle{definition}
\newtheorem*{cor*}{Corolário}
\newenvironment{corolario*}{
\definecolor{shadecolor}{HTML}{f59aa3}
\begin{shaded}
\begin{cor*}
}
{
\end{cor*}
\end{shaded}
}



\theoremstyle{definition}
\newtheorem{definition}{Definição}
\newenvironment{definicao}{
\definecolor{shadecolor}{HTML}{defcfc}
\begin{shaded}
\begin{definition}
}
{
\end{definition}
\end{shaded}
}

\theoremstyle{definition}
\newtheorem*{definition*}{Definição}
\newenvironment{definicao*}{
\definecolor{shadecolor}{HTML}{defcfc}
\begin{shaded}
\begin{definition*}
}
{
\end{definition*}
\end{shaded}
}

\theoremstyle{definition}
\newtheorem{example}{Exemplo}
\newenvironment{exemplo}{
\definecolor{shadecolor}{HTML}{b9e6d3}
\begin{shaded}
\begin{example}
}
{
\end{example}
\end{shaded}
}

\theoremstyle{definition}
\newtheorem*{example*}{Exemplo}
\newenvironment{exemplo*}{
\definecolor{shadecolor}{HTML}{b9e6d3}
\begin{shaded}
\begin{example*}
}
{
\end{example*}
\end{shaded}
}

\theoremstyle{definition}
\newtheorem{veronica}{Lema}
\newenvironment{lema}{
\definecolor{shadecolor}{HTML}{FFF1F2}
\begin{shaded}
\begin{veronica}
}
{
\end{veronica}
\end{shaded}
}

\theoremstyle{definition}
\newtheorem*{veronica*}{Lema}
\newenvironment{lema*}{
\definecolor{shadecolor}{HTML}{FFF1F2}
\begin{shaded}
\begin{veronica*}
}
{
\end{veronica*}
\end{shaded}
}

\newenvironment{prova}{\paragraph{Demonstração:}}{\hfill$\square$ 
\pulinho
}

\newenvironment{aviso}{
\definecolor{shadecolor}{HTML}{fffe9a}
\begin{shaded}
}{\end{shaded}}

\theoremstyle{definition}
\newtheorem{exersisio}{Exercício}
\newenvironment{exercicio}{
\definecolor{shadecolor}{HTML}{CFB5FF}
\begin{shaded}
\begin{exersisio}
}
{
\end{exersisio}
\end{shaded}
}

%Tamanho do papel
\usepackage[a4paper,top=1cm,bottom=1cm,margin=3.5cm]{geometry}



%%%%%%%%%%%%%%%%%%%%%%%%%%%%%%%%%%%%%%%%%%%%%%%


%Definições só pra esse texto
\usepackage[bb=boondox]{mathalpha}

\title{Pergunta do Aurichi}
\author{Violeta Mar}
\date{}

\begin{document}
\maketitle

\begin{proposicao*}
Toda extremidade de um grafo localmente finito admite um representante geodésico.
\end{proposicao*}

\begin{prova}
  Começamos argumentando que qualquer grafo localmente finito $G$ deve conter uma geodésica infinita. Para isso, fixe um vértice $v$ qualquer de $G$. Chame de $B(v,n)$ o conjunto de vértices a uma distância menor ou igual a $n \in \mathbb{N}$. Como $G$ é localmente finito, cada $B(v,n)$ é finito. Queremos juntar todos os possíveis caminhos de $G$ que começam em $v$ em um espaço, incluindo ambos os finitos e infinitos. Crie um símbolo arbitrário $\text{``fim ''}$ e defina o conjunto $B_n \defeq B(v,n) \cup \text{``fim''}$. Colocando a topologia discreta em cada um dos $B_n$, teremos pelo Teorema de Tychonoff que o produto $\prod_{n \in \mathbb{N}}^{} B_n$ é compacto. O espaço de sequências que realmente definem caminhos será chamado de $C \subset \prod_{n \in \mathbb{N}}^{} B_n$. Os elementos de $C$ são sequências que: caso apareça um ``fim'', todos os próximos elementos também são ``fim''; se temos dois elementos adjacentes $x_n$ e $x_{n+1}$, então existe uma aresta entre eles. Dentre os caminhos de $C$, defina $M \subset C$ como sendo os geodésicos (isto é, os caminhos tais que a distância entre $x_1$ e $x_n$ é sempre $n$, pra qualquer $n$ natural). É rotineiro verificar que os complementos de $C$ e de $M$ são abertos, e portanto ambos são conjuntos fechados de um compacto, e portanto são compactos.
  
  Como $G$ é infinito, cada $B(v,n)$ deve ser não vazio, e portanto deve existir uma geodésica saindo de $v$ e de tamanho $n$ para cada $n \in \mathbb{N}$. Chame esta geodésica de $\gamma_n$. Temos uma sequência infinita em $M$ compacto, e portanto ela admite um ponto de acumulação $\gamma$ em $M$. Como todo aberto ao redor de $\gamma$ contém infinitos $\gamma_n$, o $\gamma$ deve ser um caminho infinito. Assim, $\gamma$ será nossa geodésica infinita procurada.

  (sugestão do Gustavo: trocar a demonstração anterior por uma muito mais curta usando Lema de Konig)

  Vamos voltar ao objetivo original: encontrar um representante geodésico para uma dada extremidade $\epsilon$ de um grafo localmente finito $G$. Tome uma árvore geradora normal $T \subset G$ que exiba todas as extremidades de $G$, e seja $\sigma$ o raio representante em $T$ de $\epsilon$. Chame de $\tilde{G}$ o subgrafo induzido de $\sigma$ em $G$. Pelo que provamos ao início, $\tilde{G}$ deve conter uma geodésica infinita $\gamma$, que representa a mesma extremidade $\epsilon$. A normalidade de $T$ garante que $\gamma$ ser geodésica em $\tilde{G}$ implica em $\gamma$ ser geodésica em $G$, concluindo nossa demonstração.
\end{prova}





\end{document}
